\Needspace{29\baselineskip}
\subsection{Confrontation de la constante de Planck réduite}
\label{sec:revision-planck-contraste}

Le Système international en vigueur depuis le 20 mai 2019 fixe
exactement \(h=6.62607015\times10^{-34}\,\mathrm{J\,s}\)
\cite{bipm2018}. Sa constante réduite est donc
\[
 \hbar_{\mathrm{SI}}=\frac{6.62607015}{2\pi}\,10^{-34}\,
 \mathrm{J\,s}
 =1.0545718176461563912\ldots\times10^{-34}\,\mathrm{J\,s}.
\]
Les points de suspension correspondent au développement d'une expression
exacte, et non à une incertitude expérimentale de \(h\) dans le SI actuel.
Cette valeur de référence est utilisée ici après l'évaluation des
lecteurs de HMT ; sa fonction est comparative.

Les deux coordonnées du retour de l'article sont
\begin{align*}
 H_5&=1.0545718176461544696\ldots,\\
 \eta_{\mathrm{ret}}&=1.0545718169150390611\ldots.
\end{align*}
Dans l'identification de \(\mathcal U_S\) à \(\mathrm{J\,s}\), la
différence relative de la section antérieure par rapport à
\(\hbar_{\mathrm{SI}}\) est d'environ \(-1.82\times10^{-15}\).
Celle de la section postérieure est d'environ \(-6.93\times10^{-10}\).
La seconde différence contient l'effet du retour utilisé par le
correcteur électronique ; elle n'équivaut pas à l'erreur d'évaluation de la
première section.

La précision de ces comparaisons permet de reconnaître la proximité
quantitative de la construction de \(H_5\) et, simultanément, de distinguer
proximité et identité exacte. Dans le présent manuscrit, les expressions
présentées ne fournissent pas une égalité exacte avec \(h/(2\pi)\).
Le résultat démontré sur \(\Sigma_H\) est sa décade \(-34\) ;
le résultat analytique sur \(H_5\) est l'évaluation du caractère
déclaré. L'identification physique est examinée dans cette coordinatisation et avec les
différences explicites que nous venons de donner. La valeur SI n'intervient
ni dans la définition de la norme déterminantale ni dans la récurrence régionale.
